\chapter{Persiapan Lingkungan Praktikum}

\section{Perangkat yang Digunakan}

Praktikum memakai PostgreSQL 17 sebagai inti gudang data, DBeaver sebagai klien
SQL, Python 3.11 atau lebih baru untuk pipeline ETL, serta Git untuk versi dan
pengumpulan. Seluruh layanan dijalankan dari Docker sehingga satu kelas berada
pada versi yang sama.

\begin{center}
\small
\begin{tabular}{@{}p{0.26\textwidth}p{0.28\textwidth}p{0.36\textwidth}@{}}
\toprule
\textbf{Perkakas} & \textbf{Peran} & \textbf{Dipakai pada} \\
\midrule
PostgreSQL 17 & Inti gudang data & Seluruh modul \\
DBeaver & Klien SQL dan diagram E/R & Seluruh modul \\
Python 3.11+ & ETL dan pembangkitan data & Modul 7, proyek akhir \\
Git & Versi dan pengumpulan & Seluruh modul \\
Docker Compose & Reproduksibilitas lingkungan & Seluruh modul \\
dbt Core & Transformasi, dokumentasi, lineage & Modul 8, proyek akhir \\
Apache Superset & Dashboard dan eksplorasi visual & Modul 9, proyek akhir \\
\bottomrule
\end{tabular}
\end{center}

Pustaka Python yang diperlukan: \perintah{psycopg2-binary}, \perintah{SQLAlchemy},
\perintah{pandas}, dan \perintah{Faker}. Seluruhnya tercantum pada
\berkas{requirements.txt}.

dbt Core dan Superset bersifat pengenalan. Keduanya ditunjukkan dan dipakai
secara terbimbing, tetapi tidak dinilai sedalam PostgreSQL dan Python.

\section{Instalasi}

Mahasiswa memasang Docker Desktop, DBeaver Community, Git, dan Python 3.11+ di
laptop masing-masing. Sisanya disediakan oleh \berkas{docker-compose.yml}.

\begin{lstlisting}[style=shellgaya]
git clone <url-repositori-kelas> praktikum-pgd
cd praktikum-pgd
cp etl/config/.env.example .env
docker compose pull          # dikerjakan di rumah, bukan di laboratorium
docker compose up -d
docker compose ps
\end{lstlisting}

Lingkungan menyediakan tiga layanan. Pemisahan basis data sumber dan gudang
sejak awal membuat batas keduanya terasa nyata: tidak ada query yang dapat
menyeberang begitu saja dari satu ke yang lain.

\begin{center}
\small
\begin{tabular}{@{}p{0.24\textwidth}p{0.10\textwidth}p{0.26\textwidth}p{0.30\textwidth}@{}}
\toprule
\textbf{Layanan} & \textbf{Port} & \textbf{Basis data} & \textbf{Peran} \\
\midrule
\perintah{postgres\_source} & 5433 & \perintah{nusamart\_oltp} &
  Sistem operasional NusaMart \\
\perintah{postgres\_dw} & 5434 & \perintah{nusamart\_dw} &
  Gudang data yang dibangun mahasiswa \\
\perintah{superset} & 8088 & --- & Dashboard, mulai Modul 9 \\
\bottomrule
\end{tabular}
\end{center}

Port sengaja digeser dari 5432 agar tidak berbenturan dengan PostgreSQL yang
mungkin sudah terpasang langsung di laptop.

\subsection{Ekstensi PostgreSQL}

Ekstensi berikut diaktifkan asisten pada basis data gudang sebelum modul yang
memerlukannya.

\begin{table}[htbp]
\centering
\small
\caption{Ekstensi PostgreSQL yang perlu diaktifkan}
\label{tab:ekstensi}
\begin{tabular}{@{}p{0.24\textwidth}p{0.54\textwidth}p{0.12\textwidth}@{}}
\toprule
\textbf{Ekstensi} & \textbf{Untuk apa} & \textbf{Modul} \\
\midrule
\perintah{btree\_gist} & Exclusion constraint penjaga versi SCD-2 & 3 \\
\perintah{postgres\_fdw} & Membaca basis data sumber lain secara langsung & 7 \\
\perintah{tablefunc} & \perintah{crosstab} untuk pivot, opsional & 9 \\
\perintah{pg\_stat\_statements} & Pemantauan query operasional & Proyek akhir \\
\bottomrule
\end{tabular}
\end{table}

\section{Uji Lingkungan}

Jalankan pemeriksaan berikut sebelum memulai setiap modul. Ketiganya harus
memberi keluaran, bukan pesan galat.

\begin{lstlisting}[style=shellgaya]
docker compose ps

docker compose exec postgres_source \
  psql -U praktikum -d nusamart_oltp -c "SELECT version();"

docker compose exec postgres_dw \
  psql -U praktikum -d nusamart_dw   -c "SELECT current_database();"
\end{lstlisting}

Untuk modul yang memakai Python, tambahkan pemeriksaan berikut.

\begin{lstlisting}[style=pythongaya]
import sys, psycopg2, sqlalchemy, pandas as pd

print("Python     :", sys.version.split()[0])
print("psycopg2   :", psycopg2.__version__)
print("SQLAlchemy :", sqlalchemy.__version__)
print("pandas     :", pd.__version__)
\end{lstlisting}

\section{Aturan Reproduksibilitas}

Setiap luaran wajib dapat dijalankan ulang dari awal oleh orang lain di laptop
yang berbeda. Karena itu berlaku aturan berikut sepanjang praktikum.

\begin{enumerate}[leftmargin=1.4em]
  \item Skrip DDL bersifat \istilah{idempoten}: memakai
        \perintah{DROP ... IF EXISTS} atau \perintah{CREATE ... IF NOT EXISTS}
        sehingga dapat dijalankan dua kali tanpa galat.
  \item Pipeline ETL harus idempoten: dijalankan dua kali berturut-turut, jumlah
        baris tidak berubah dan tidak ada versi dimensi yang kembar.
  \item Kredensial tidak pernah ditulis di dalam skrip. Nilainya dibaca dari
        \berkas{.env}, dan \berkas{.env} tidak masuk version control.
  \item Nilai \perintah{--seed} pembangkit data dicatat pada laporan. Tanpa
        seed, angka hasil profiling tidak dapat dibandingkan.
  \item Path ditulis relatif terhadap akar repositori, bukan path absolut yang
        hanya ada di laptop sendiri.
  \item Data mentah bersifat \istilah{read-only}. Perbaikan kualitas data
        dilakukan pada \perintah{staging} atau \perintah{gudang}, tidak pernah
        pada basis data sumber.
\end{enumerate}

\section{Penyiapan Data}

Data NusaMart dibangkitkan dengan \perintah{Faker} oleh skrip yang disediakan
asisten, dan tersedia dalam tiga ukuran. Penjenjangan ini penting: pada 100 ribu
baris, \istilah{sequential scan} justru menang, sehingga Modul 6 tidak akan
memperlihatkan apa pun bila memakai ukuran kecil.

\begin{center}
\small
\begin{tabular}{@{}p{0.16\textwidth}p{0.26\textwidth}p{0.48\textwidth}@{}}
\toprule
\textbf{Ukuran} & \textbf{Baris fakta} & \textbf{Dipakai pada} \\
\midrule
Kecil  & $\approx$100 ribu & Modul 2--5 \\
Sedang & $\approx$5 juta   & Modul 6--8, wajib pada Modul 6 \\
Besar  & $\approx$50 juta  & Opsional, bagi yang ingin melihat partisi dan
                             agregat benar-benar berbeda \\
\bottomrule
\end{tabular}
\end{center}

Mulai Modul 7, dua sumber operasional tambahan disediakan dengan
ketidaksepakatan yang sengaja ditanam: penamaan kolom berbeda, satuan berat
berbeda antara gram dan kilogram, kode kategori berbeda, pelanggan yang sama
tercatat dengan ejaan berbeda, dan sebagian transaksi datang terlambat tiga
hari. Basis data Northwind disediakan terpisah sebagai pembanding pada Modul 6.

Artefak besar dan data yang dapat dibangkitkan ulang tidak disimpan ke
repositori. Data asli diletakkan sebagai read-only di \berkas{data/raw/} dan
hasil proses di \berkas{data/processed/}.
